\subsection{不变微分形式的构造}

引理 2. 设 \(\mathbf{G}\) 是李群，\(\mathrm{U}\) 是 \(e\) 在 \(\mathbf{G}\) 中的一个对称开邻域，\(\mathrm{E}\) 是完备可赋范空间，且 \(\phi: \mathrm{U}^2 \to \mathrm{E}\) 是解析映射。对任意 \(g \in \mathrm{U}\)，设 \(\omega_g\) 是某映射在点 \(g\) 处的微分，该映射为 \(h \mapsto \phi(g^{-1}h)\)。则 \(\omega\) 是 \(\mathrm{U}\) 上的左不变微分形式，它是 \(\mathbf{G}\) 上在 \(e\) 处的值为 \(d_e\phi\) 的左不变微分形式的限制。

显然 \(\omega_{e} = d_{e}\phi\)。对任意 \(g \in \mathrm{U}\) 及 \(t \in \mathrm{T}_{e}(\mathrm{G})\)，

\[
\begin{array}{r l} \langle \omega_ {g}, \mathrm{T} _ {e} (\gamma (g)) t \rangle & = \langle d _ {g} (\phi \circ \gamma (g) ^ {- 1}), \mathrm{T} _ {e} (\gamma (g)) t \rangle \\ & = \langle d _ {e} \phi \circ \mathrm{T} _ {g} (\gamma (g) ^ {- 1}), \mathrm{T} _ {e} (\gamma (g)) t \rangle = \langle d _ {e} \phi , t \rangle \end{array}
\]

因此 \(\omega_{g}\) 由 \(d_{e}\phi\) 经 \(\mathrm{T}_{e}(\gamma(g))\) 导出。

命题 52. 设 \(n\) 是整数 \(>0\)，\(G\) 是 \(n\) 维李群，\(U\) 是 \(e\) 在 \(G\) 中的对称开邻域，且 \(\psi: U^2 \to K^n\) 是 \(G\) 的一个图，满足 \(\psi(e) = 0\)。若 \((x_1, \ldots, x_n)\) 是 \(x \in \psi(U)\) 的坐标，\((y_1, \ldots, y_n)\) 是 \(y \in \psi(U)\) 的坐标，我们记

\[
m _ {1} (x _ {1}, \dots , x _ {n}, y _ {1}, \dots , y _ {n}), \dots , m _ {n} (x _ {1}, \dots , x _ {n}, y _ {1}, \dots , y _ {n})
\]

\(\psi (\psi^{-1}(x)^{-1}\psi^{-1}(y))\) 的坐标。于是，若对 \(1\leqslant k\leqslant n\) 写成

\[
\begin{array}{r l} \varpi_ {k} (x _ {1}, \ldots , x _ {n}) = D _ {n + 1} m _ {k} (x _ {1}, \ldots , x _ {n}, x _ {1}, \ldots , x _ {n}) d x _ {1} + \dots \\ & \quad + \mathrm{D} _ {2 n} m _ {k} (x _ {1}, \ldots , x _ {n}, x _ {1}, \ldots , x _ {n}) d x _ {n}, \end{array}\tag{32}
\]

则微分形式 \(\varpi_k\) 定义在 \(\psi(\mathbf{U})\) 上，由 \(\psi\) 从 \(\mathbf{G}\) 上的左不变微分形式导出，并且满足 \(\varpi_k(0,\ldots,0) = dx_k\)。

应用引理2，取 E = K，并令 \(\phi(g)\) 为 \(\psi(g)\) 的指标为 k 的坐标。得到微分形式 \(\omega_{k}\)；设 \(\varpi_{k}\) 是它在 \(\psi\) 下的变换。\(\varpi_{k}\) 在 \((x_{1},\ldots,x_{n})\) 处的值，是在 \((x_{1},\ldots,x_{n})\) 处对函数 \(y\mapsto m_{k}(x_{1},\ldots,x_{n},y_{1},\ldots,y_{n})\) 求得的微分；因此该值由公式 (32) 给出。于是只需使用引理2的结论。

命题 53. 设 G 是李群，A 是完备可赋范代数，且 \(\phi\) 是从 G 到 A* 的李群态射。对任意 \(g \in G\)，设 \(\omega_g = \phi(g)^{-1}.d_g\phi\)。则 \(\omega\) 是 G 上在 \(e\) 处取值为 \(d_e\phi\) 的左不变微分形式。

应用引理2，取 \(\mathbf{E} = \mathbf{A}\) 且 \(\mathbf{U} = \mathbf{G}\)。在 \(g\) 处，映射 \(h \mapsto \phi(g^{-1}h) = \phi(g)^{-1}\phi(h)\) 的微分是 \(\phi(g)^{-1}.d_g\phi\)。

